\documentclass[11pt]{article}
% Preámbulo compartido para los modelos de parcial ACN
\usepackage[utf8]{inputenc}
\usepackage[T1]{fontenc}
\usepackage[spanish]{babel}
\usepackage{amsmath,amssymb}
\usepackage[margin=2.3cm]{geometry}
\usepackage{enumitem}
\usepackage{fancyhdr}
\usepackage{array}

\pagestyle{fancy}
\fancyhf{}
\lhead{Aplicaciones Computacionales en Negocios --- UTDT}
\rhead{\thepage}
\renewcommand{\headrulewidth}{0.4pt}

% Encabezado tipo examen
\newcommand{\encabezado}[1]{%
  \begin{center}
    {\Large\bfseries Examen parcial}\\[2pt]
    {\large Aplicaciones Computacionales en Negocios --- UTDT}\\[2pt]
    {\normalsize #1}\\[8pt]
  \end{center}
  \noindent\fbox{\parbox{\linewidth}{\centering\bfseries RESPONDER EXCLUSIVAMENTE DEBAJO DEL ESPACIO PROVISTO BAJO CADA PREGUNTA}}
  \vspace{6pt}

  \noindent Nombre y apellido: \underline{\hspace{7cm}} \hfill Legajo: \underline{\hspace{3.5cm}}
  \vspace{6pt}
  \hrule
  \vspace{10pt}
}

% Puntaje en negrita
\newcommand{\pts}[1]{\textbf{[#1 puntos]}\ }

% Espacio de respuesta
\newcommand{\espacio}[1]{\vspace{#1}}


\begin{document}
\encabezado{Modelo 3}

% ------------------------------------------------------------------
\noindent\textbf{1)} Una empresa opera dos parques eólicos, Norte y Sur. La velocidad regional del viento en el día $i$ sigue el modelo con reversión a la media estacional
\[
S_{i+1} = S_i + \kappa\,(\theta(t_i) - S_i)\,\Delta + \sigma\sqrt{\Delta}\,Z_i, \qquad \theta(t_i) = 6 + 2\cos(2\pi i / 365)
\]
con $\Delta = 1/365$ y $Z_i$ shocks normales estándar independientes entre días.

\medskip
\pts{10} ¿Qué rol cumple cada parámetro ($\kappa$, $\theta$, $\sigma$)? ¿Cómo es el gráfico de la evolución temporal de $S$ si $\sigma = 0$? Dibujalo esquemáticamente para un año.

\espacio{4.5cm}

\pts{15} La diferencia de viento entre ambos molinos arranca en cero y evoluciona como
\[
D_{i+1} = D_i - \beta\,D_i\,\Delta + \gamma\sqrt{\Delta}\,W_i, \qquad D_0 = 0
\]
con $W_i$ shocks normales estándar, independientes de los $Z_i$. Explicá el comportamiento de $D$ para distintos valores de $\beta$. ¿Cuál es su distribución de largo plazo?

\espacio{4.5cm}

% ------------------------------------------------------------------
\noindent\textbf{2)} Suponé que tenés disponible una función para generar variables uniformes $[0,1]$ y otra para generar variables normales estándar.

\medskip
\pts{25} Describí de manera precisa (pseudocódigo y palabras) un algoritmo para simular \emph{simultáneamente} un camino de $S$ (viento regional) y un camino de $D$ (diferencia entre molinos) durante un año completo ($i = 0, \dots, 365$), a partir de las ecuaciones del punto 1. A partir de $S$ y $D$, ¿cómo obtenés la velocidad en cada molino individual?

\espacio{6.5cm}

% ------------------------------------------------------------------
\noindent\textbf{3)} Cada molino puede romperse por viento excesivo. La probabilidad de que el molino se rompa en el día $i$, si estaba activo el día anterior, es $0.25\, S_i^2\, \Delta$. Si se rompe, ese día no produce y se repara ese mismo día a un costo fijo. El ingreso diario es $1000 \cdot S_i$ dólares si el molino está activo.

\medskip
\pts{25} Describí con precisión (pseudocódigo y palabras) cómo simular, sobre un año, la evolución conjunta del viento, las roturas y la acumulación de cashflows de un molino. Explicá cómo incorporás la decisión de rotura día a día usando la función de uniformes.

\espacio{6.5cm}

\pts{15} A partir de un número grande de caminos simulados, describí cómo estimar el ingreso neto anual esperado del molino y su error de estimación. Si querés reducir ese error a la décima parte, ¿cuántos caminos más necesitás? Justificá.

\espacio{4.5cm}

\end{document}
